\section{Comparison of the properties of morphisms}
\label{XII.3}

\begin{proposition}
\label{XII.3.1} Let $f \colon X \to Y$ be a morphism of
$\CC$-schemes locally of finite type, $f^\an \colon X^\an \to
Y^\an$ the morphism deduced from~$f$ on the associated
analytic spaces. Let $P$ be the property of being
\begin{enumerate}
\item[(i)] flat
\item[(ii)] net (\ie unramified)
\item[(iii)] étale
\item[(iv)] smooth
\item[(v)] normal
\item[(vi)]
% original p. 321
reduced
\item[(vii)] injective
\item[(viii)] separated
\item[(ix)] an isomorphism
\item[(x)] a monomorphism
\item[(xi)] an open immersion.
\end{enumerate}
Then, for $f$ to have the property $P$, it is necessary and
sufficient that $f^\an$ have it.
\end{proposition}

Let us denote by $\varphi \colon X^\an \to X $ and $\psi \colon Y^\an \to Y$ the
canonical morphisms. Let $x$ be a point of~$X^\an$, $y = f^\an
(x)$. The morphisms $\cal{O}_{Y^\an, y} \to {\mathcal O}_{X^\an, x}$
and $\cal{O}_{Y, \psi(y)} \to \cal{O}_{X, \varphi(x)} $ deduced from
$f$ and $f^{\an}$ give the same morphism on passing to the
completions \eqref{XII.1.1}. By \cite[ch\ptbl 3 \S 5 prop\ptbl 4]{XII.2}
(\resp EGA~IV 17.4.4) it therefore amounts to the same to say that $f^\an $
satisfies property (i) (\resp (ii)) or to say that $f$
satisfies (i) (\resp (ii)) at each closed point of~$X$. As
the set of points of~$X$ where (i) (\resp (ii)) is
satisfied is an open set (EGA~IV 11.1.1 and I 3.3), this
proves (i) and (ii), hence also (iii).

Let $P$ be the property (iv) (\resp (v), \resp (vi)). Taking into account
\Ref{XII.2.1} ((v), (vi), (vii)), it amounts to the same to say that
the geometric fibers of~$f^\an$ at the various points
$y$ of~$Y^\an$ are regular (\resp normal, \resp reduced) or that the same holds for the geometric fibers of
$f$ at the various closed points $\psi (y)$ of~$Y$. The cases
(iv) (\resp (v), \resp (vi)) then follow from (i) and from the fact that
the set of points of~$Y$ where the geometric fibers of
$f$ are regular is an open set (EGA~IV 12.1.7).

(vii). If $f$ is injective, so is $f^\an$.
Conversely, suppose $f^\an$ injective and let us show that the same holds
for~$f$. We may assume
% original p. 322
$f$ of finite type. The morphism $f^\an$ being injective, the fibers
of~$f$ at the closed points of~$Y$ are radicial; as the set
of points of~$Y$ whose fiber is radicial is locally
constructible (EGA~IV 9.6.1) and as $Y$ is a Jacobson
scheme, $f$ has all its fibers radicial, hence is injective.

(viii). Let $\Delta \colon X \to X \times_Y X $ and $\Delta^\an
\colon X^\an \to X^\an \times_{Y^\an} X^\an $ be the diagonal
immersions, $\Theta \colon X^\an \times_{Y^\an} X^\an \to X \times_Y X
$ the canonical morphism. By virtue of~\Ref{XII.2.3} it amounts to the
same to say that $\Delta (X)$ is closed in $X \times_Y X $
or that $\Delta^\an (X^\an)$ is closed in $X^\an \times_{Y^\an}
X^\an$.

As an open immersion is nothing other than an injective étale
morphism (EGA~IV 17.9.1 and \cite[\no 13 \S 1]{XII.4}), (xi) follows from (iii)
and (vii). An isomorphism being the same thing as a surjective
open immersion, (ix) follows from (xi) and from
(\Ref{XII.3.2}~(i)) below. To say that $f$ is a monomorphism
is equivalent to saying that the diagonal morphism $\Delta \colon X \to
X \times_Y X$ is an isomorphism, hence (x) follows from (ix).

\begin{proposition}
\label{XII.3.2}
Let $X$ and $Y$ be two $\CC$-schemes locally of finite type,
$f\colon X \to Y$ a morphism of finite type, $f^\an \colon X^\an \to
Y^\an $ the morphism deduced from~$f$ on the associated
analytic spaces. Let $P$ be the property of being
\begin{enumerate}
\item[(i)] surjective
\item[(ii)] dominant
\item[(iii)] a closed immersion
\item[(iv)] an immersion
\item[(v)] proper\footnote{We shall say that a morphism of analytic
spaces is proper if it is so in the sense of \cite[ch\ptbl 1 \S 10 \no 1]{XII.1} and
if it is separated.}
\item[(vi)]
% original p. 323
finite.
\end{enumerate}
Then, for $f$ to have the property $P$, it is necessary and
sufficient that $f^\an$ have it.
\end{proposition}

Let $\varphi \colon X^\an \to X$ and $\psi \colon Y^\an \to Y$ be the
canonical morphisms.

(i). If $f$ is surjective, for every point $y$ of~$Y^\an$, $f^{-1}
(\psi (y)) $ is a nonempty closed subset of~$X$; it therefore
contains at least one closed point, which proves that $f^\an$ is
surjective. Conversely, if $f^\an$ is surjective, $f(X)$ is a
locally constructible subset of~$Y$ (EGA~IV 1.8.4) which contains
all the closed points of~$Y$; we therefore have $f(X)=Y$.

(ii) follows from~\Ref{XII.2.2}.

(iii). If $f$ is a closed immersion, so is
$f^\an$ by \Ref{XII.1.1}~a). Conversely, if $f^\an$ is a closed
immersion, so is~$f$ by
\Ref{XII.3.1}~(x) and \Ref{XII.3.2}~(v), since this amounts to saying
that $f$ is a proper monomorphism (EGA~IV 8.11.5).

(iv). It is clear that, if $f$ is an immersion, so is~$f^\an$.
Conversely, suppose that $f^\an$ is an
immersion, and let $T$ be the image of~$X$ in $Y$, $\overline{T}$
the scheme-theoretic closure of~$f$. We have a factorization of~$f$,
$$
X \lto{i} \overline{T} \lto{j} J \quoi,
$$
where $j$ is a closed immersion, $i$ the canonical morphism,
and we deduce from it the following factorization of~$f^\an$
$$
X^\an \lto{i^\an} \overline{T}^{\mkern1mu \an} \lto{j^\an} Y^\an \
.
$$
As
% original p. 324
$T = f(X)$ is a locally constructible subset of~$Y$ (EGA~IV
1.8.4), we have, by~\Ref{XII.2.2}, $\overline{T}^{\mkern1mu \an} =
\overline{f^\an (X^\an )}$. It follows that $i^\an (X^\an )$ is
an open subset of~$\overline{T}^{\mkern1mu \an}$, hence that $i(X) $ is an open subset of
$\overline{T}$. Consider the canonical factorization of~$i$
$$
X \lto{i_1} i (X) \lto{i_2} \overline{T}.
$$
The morphism $i_1^\an $ is a proper monomorphism, hence so is
$i_1$ by \Ref{XII.3.2}~(v) and \Ref{XII.3.1}~(x);
this proves that $i_1$, hence also $f$, is an immersion.

(v). Suppose that $f$ is proper and let us show that so is~$f^\an$.
Since the fact that $f^\an $ is proper is local
over~$Y^\an$, we may assume $Y$ affine. By Chow's lemma
(EGA~II 5.6.1), we can find a projective $Y$-scheme $X'$ and a
projective surjective morphism
$$
g \colon X' \to X.
$$
The morphism $(fg)^\an = f^\an g^\an $ is projective hence proper,
$g^\an $ is surjective, and it follows from \cite[ch\ptbl 1 \S 10]{XII.1} that
$f^\an $ is proper.

Conversely, suppose $f^\an $ proper and let us show that so is~$f$.
By \Ref{XII.3.1}~(viii) $f$ is
separated. It remains to prove that $f$ is universally
closed, and it even suffices to show that $f$ is closed; indeed,
for every $Y$-scheme $Y'$ locally of finite type, the
morphism
$$
f_{(Y')} = h \colon X \times_Y Y' \to Y'
$$
will also be closed since $h^\an $ is proper. Let $T$ be a closed
subset of~$X$; $f(T)$ is a locally constructible set, and
we have
$$
f^\an (\varphi^{-1} (T)) = \psi^{-1} (f(T)).
$$
As
% original p. 325
$f^\an $ is proper, $\psi^{-1} (f (T))$ is a closed subset of
$Y^\an $, and it therefore follows from~\Ref{XII.2.2} that we have
$$
\psi^{-1} (\overline{f(T)}) = \psi^{-1} (f(T))\quoi .
$$
This implies that we have $f(T) = \overline{f(T)}$, \ie that $f$
is closed, hence that $f$ is proper.

(vi). It amounts to the same to say that a morphism is finite or that it
is proper with finite fibers (EGA~III 4.4.2 and \cite[\no 19 \S 5]{XII.4}).
As the set of points where the fibers of~$f$ are finite is
locally constructible (EGA~IV 9.7.9), the fibers of~$f$ are finite
if and only if the same holds for the fibers of~$f^\an $; (vi)
therefore follows from (v).

\begin{remark}
\label{XII.3.3}
\begin{enumerate}
\item[a)] let $f \colon X \to Y$ be a morphism of~$\CC$-schemes
locally of finite type. The fact that $f^\an $ is a local
isomorphism does not imply that the same holds for~$f$. Indeed,
if $f$ is étale, $f^\an $ is étale hence is a local
isomorphism \hbox{\cite[\no 13 \S 1]{XII.4}}, but this is not necessarily the case
for~$f$.
\item[b)]
Statement~\Ref{XII.3.2}
is not true if one does not assume
$f$ of finite type. Let us show for example that $f^\an $ can be
a closed immersion without $f$ being one. Indeed, it
suffices to take for $X$ the sum of~$\ZZ $ copies of
$\Spec \CC$, and for $Y$ the affine line, and for $f$ the morphism
obtained by sending the points of~$X$ to distinct points of~$Y$
forming a discrete subset.
\end{enumerate}
\end{remark}

\section[Cohomological comparison theorems]{Cohomological comparison theorems and existence theorems}
\label{XII.4}
% original p. 326

The object of this no.\ is to prove anew the results of
\cite[\no 2 th\ptbl 5 and th\ptbl 6]{XII.3}; the latter generalize to the case
of a proper scheme the theorems established in \cite[\no 12]{XII.10} when $X$ is projective, and extend them to the
relative case. More general results, concerning
relative schemes proper over an analytic space, are
proved in \cite[ch\ptbl VIII \no 3]{XII.7}.

Recall that the {\v{C}}ech cohomology used in \cite[\no 12]{XII.10} coincides with the usual cohomology in the
algebraic case as in the analytic case (EGA~III 1.4.1. and \cite[II 5.10]{XII.5}).
\subsection{}
\label{XII.4.1}
Let $f \colon X \to Y $ be a morphism of~$\CC$-schemes
locally of finite type and consider the commutative diagram
$$
\xymatrix{ X^\an \ar[r]^\varphi \ar[d]_{f^\an} & X \ar[d]^f
\\ Y^\an \ar[r]^\psi & Y \quoi.}
$$
If $F$ is an $\cal{O}_X$-Module, we have, for every integer $p \geq
0$, morphisms
$$
\R^p f_* F \lto{i} \R^p f_* (\varphi_* F^\an )
\lto{j} \R^p (f\cdot\varphi)_* F^\an \lto{k} \psi_*
(\R^p f_*^\an F^\an ) \quoi,
$$
where $i$ is deduced from the canonical morphism $F \to \varphi_*
F^\an $, and $j, k$ are ``edge homomorphisms'' of Leray
spectral sequences. To the composite $k.j.i$ there is associated a
canonical morphism
\begin{equation*}
\label{eq:XII.4.1.1}
\tag{4.1.1} {\theta_p \colon (\R^p f_* F)^\an \to \R^p f_*^\an
(F^\an )}
\end{equation*}

\begin{theorem}
\label{XII.4.2}
Let
% original p. 327
$f \colon X \to Y $ be a proper morphism of~$\CC$-schemes
locally of finite type, $F$ a coherent
$\cal{O}_X$-Module. Then, for every integer $p \geq 0$ the morphism
\eqref{eq:XII.4.1.1}
$$
\theta_p \colon (\R^p f_* F)^\an \to \R^p f_*^\an (F^\an)
$$
is an isomorphism.
\end{theorem}

1) \emph{Case where $f$ is projective}. The proof is
analogous to that of \cite[\no 13]{XII.10}. Let us recall it briefly. One
reduces to the case where $X$ is a standard projective space $\PP^r_Y
$ over~$Y$. Let $\cal{Y} = Y^\an$, $\cal{P} = \PP^r_\cal{Y}$;
one first proves that we have
$$
f_*^\an \cal{O}_\cal{P} = \cal{O}_\cal{Y} \quoi, \qquad \R^p f_*^\an
(\cal{O}_\cal{P}) = 0 \qquad \text{for } p>0
$$
To verify the preceding relations, one can indeed
reduce to the case where $\cal{Y}$ is a ball $\cal{B}$ of an
affine space $\cal{E}^n$. One considers the ``standard covering'' $\{
\cal{U}_i \}$ of~$\cal{P}$ by $r+1$ open sets isomorphic to $\cal{B}
\times \cal{E}^r$. As these open sets are Stein, we have, for every
integer $p \geq 0$, isomorphisms
$$
\H^p (\{ \cal{U}_i \}, \cal{O}_\cal{P}) \isomto \H^p (\cal{P},
\cal{O}_\cal{P})
$$
One can then express the sections of the structure sheaf
$\cal{O}_\cal{P}$ over the open sets $\cal{U}_i $ and over their
intersections in terms of Laurent series; an easy computation
proves that we have
$$
\H^0 (\cal{P}, \cal{O}_\cal{P} ) \isomto \H^0 (\cal{Y},
\cal{O}_\cal{Y} ) \quoi, \qquad \H^p (\cal{P}, \cal{O}_\cal{P} ) = 0
\qquad \text{for } p>0
$$
The proof is then completed by copying \cite[\no 12
lemma~5]{XII.10}, the cohomology groups being replaced by the
cohomology sheaves.

2)
% original p. 328
\emph{Case where $f$ is proper}. One uses EGA~III 3.1.2 to
reduce to the projective case. Let $\cal{K}$ be the category of
coherent $\cal{O}_X$-Modules such that $\theta_p$ is an
isomorphism for every $p \geq 0$. It suffices to prove that, for
every exact sequence $0 \to F' \to F \to F'' \to 0$
of which two terms are in $\cal{K}$, the same holds for the
third, that a direct factor of an object of~$\cal{K}$ is in
$\cal{K}$, and that, for every point $x$ of~$X$, one can find an
object $F$ of~$\cal{K}$ such that we have $F_x \neq 0$.

The first condition follows by application of the five
lemma to the following commutative diagram, whose rows are exact:
$$
\xymatrix@=.7cm{ \ar[r] & (\R^p f_* F')^\an \ar[d] \ar[r] &(\R^p f_*
F)^\an \ar[d] \ar[r] & (\R^p f_* F'')^\an \ar[d] \ar[r] &
(\R^{p+1} f_* F' )^\an \ar[r] \ar[d] &\, \\ \ar[r] & \R^p
f_*^\an F'^\an \ar[r] &\R^p f_*^\an F^\an \ar[r] & \R^p
f_*^\an F''^\an \ar[r] & \R^{p+1} f_*^\an F'^\an \ar[r] &
\quoi,}
$$
and one verifies the second condition in an analogous way.

To verify the third condition, one can restrict oneself to the case
where $X$ is an irreducible scheme with generic
point $x$. We could have assumed $Y$ noetherian from
the start. By Chow's lemma (EGA~II 5.6.1), one can
find a projective $Y$-scheme~$X'$ and a projective
surjective morphism $g \colon X' \to X$. On the other hand there exists an integer~$n$
such that we have $\R^p g_* (\cal{O}_{X'} (n)) = 0 $ for every $p>0$ and
such that the canonical morphism $g^* g_* (\cal{O}_{X'} (n)) \to
\cal{O}_{X'} (n)$ is surjective (EGA~III 2.2.1). If we set
$F = g_* (\cal{O}_{X'} (n)) $, the sheaf $F$ answers
the question. Indeed, we have $F_x \neq 0$; moreover, the Leray
spectral sequence
$$
\R^p f_* (\R^q g_* (\cal{O}_{X'} (n))) \To \R^{p+q}
(f\cdot g)_* (\cal{O}_{X'} (n))
$$
being degenerate, we have an isomorphism
$$
\R^p f_* F \isomto \R^p (f\cdot g)_* (\cal{O}_{X'} (n))
$$
% original p. 329
As in the algebraic case, we have a canonical isomorphism
$$\R^p f_*^\an F^\an \isomto
\R^p (f\cdot g)_*^\an (\cal{O}_{X'}(n)^\an ),$$
and the diagram
$$
\xymatrix@C=1cm{
(\R^p f_* F)^\an \ar[d]_{\theta_p} \ar[r]^-{\sim} & (\R^p (f\cdot g)_* (\cal{O}_{X'} (n)))^\an \ar[d]_{\psi_p} \\
\R^p f_*^\an F^\an \ar[r]^-{\sim} & \R^p (f\cdot g)_*^\an
(\cal{O}_{X'}(n)^\an)}
$$
is commutative. By~1) $\psi_p $ is an isomorphism; hence so
is~$\theta_p$, which completes the
proof.

\begin{corollary}
\label{XII.4.3}
Let $X$ be a proper $\CC$-scheme, $F$ a coherent $\cal{O}_X$-Module.
Then, for every integer $p\geq 0$, the canonical
morphism
$$
\H^p(X,F)\to \H^p(X^{\an},F^{\an})
$$
is an isomorphism.
\end{corollary}

\begin{theorem}
\label{XII.4.4}
Let $X$ be a proper $\CC$-scheme. The functor which associates to every
coherent $\cal{O}_X$-Module $F$ its inverse image
$F^{\an}$ on~$X^{\an}$ is an equivalence of categories.
\end{theorem}

1) \emph{The functor is fully faithful}. Indeed, let $F$
and $G$ be two coherent $\cal{O}_X$-Modules. The canonical morphism
$$
\Hom_{\cal{O}_X}(F,G)\to \Hom_{\cal{O}_{X^{\an}}}(F^{\an},G^{\an})
$$
is identified with the canonical morphism
$$
\H^0(X,\SheafHom_{\cal{O}_X}(F,G))\to
\H^0(X^{\an},\SheafHom_{\cal{O}_X}(F,G))
$$
% original p. 330
(EGA~$0_{\textup{I}}$~6.7.6). As $\SheafHom_{\cal{O}_X}(F,G)$ is
coherent, it follows from~\Ref{XII.4.3} that this morphism is
bijective.

2) \emph{The functor is essentially surjective}. When $X$ is
projective, the assertion follows from \cite[\no 12~th\ptbl 3]{XII.10}. The
general case reduces to the preceding one by using
Chow's lemma (EGA~II~5.6.1). Indeed, let $X'$ be a projective
$\CC$-scheme, $f\colon X'\to X$ a projective surjective morphism, $U$ a
dense open subset of~$X$ such that $f$ induces an isomorphism
$f^{-1}(U)\simeq U$. We argue by noetherian induction
on~$X$; we may therefore assume that, for every coherent sheaf
$\cal{G}$ on~$X^{\an}$ such that one can find a closed
subset $Y$ of~$X$ distinct from~$X$, satisfying the relation
$Y^{\an}\supset \Supp \cal{G}$, there exists a coherent sheaf $G$
on~$X$ such that we have an isomorphism $G^{\an}\simeq \cal{G}$.

Let $\cal{F}$ be a coherent sheaf of modules on~$\cal{O}_{X^{\an}}$, and $\cal{K}$ and $\cal{L}$ the coherent
sheaves defined by the condition that the sequence
$$
0\to\cal{K}\to\cal{F}\to f_*^{\an}f^{\an *}\cal{F}\to\cal{L}\to 0
$$
be exact. As $X'$ is projective, there exists a coherent
$\cal{O}_{X'}$-Module $F'$ such that we have
$F'^{\an}\simeq {f^{\an}}^*\cal{F}$; one then deduces
from~\Ref{XII.4.2} that we have an isomorphism $(f_*F')^{\an}\simeq
f_*^{\an}{f^{\an}}^*\cal{F}$. As $\cal{K}|U^{\an}$ and
$\cal{L}|U^{\an}$ are zero, there exist coherent $\cal{O}_X$-Modules
$K$ and $L$ such that we have isomorphisms
$K^{\an}\simeq \cal{K}$, $L^{\an}\simeq \cal{L}$. By~1) the
morphism $f_*^{\an}{f^{\an}}^*\cal{F}\to\cal{L}$ comes from a unique
morphism $f_*F'\to L$; let $I=\Ker(f_*F'\to L)$. The sheaf
$\cal{F}$ is then an extension of~$I^{\an}$ by~$K^{\an}$, and it suffices
to see that this extension comes
% original p. 331
by inverse image from an extension of~$I$ by~$K$. It therefore suffices to
prove that the canonical morphism
\begin{equation*}
\label{eq:XII.4.*}
\tag{$*$} \Ext_{\cal{O}_X}^q(I,K)^{\an}\isomto
\Ext_{\cal{O}_{X^{\an}}}^q(I^{\an},K^{\an})\quad q\neq 1
\end{equation*}
is bijective. Now we have isomorphisms
$\mathbf{Ext}_{\cal{O}_X}^q(I,K)^{\an}\isomto
\mathbf{Ext}_{\cal{O}_{X^{\an}}}^q(I^{\an},K^{\an})$ for every integer
$q\geq 0$ (EGA~$0_{\textup{III}}$~12.3.5), and a morphism of spectral
sequences
$$
\xymatrix@R=.7cm@C=.5cm{
\H^p(X,\mathbf{Ext}_{\cal{O}_X}^q(I,K))\ar@{=>}[r]\ar[d] &
\Ext_{\cal{O}_X}^{p+q}(I,K)\ar[d]\\
\H^p(X^{\an},\mathbf{Ext}_{\cal{O}_{X^{\an}}}^q(I^{\an},K^{\an}))\ar@{=>}[r] & \Ext_{\cal{O}_{X^{\an}}}^{p+q}(I^{\an},K^{\an}).
}
$$
This morphism is an isomorphism since, by~\Ref{XII.4.3}, it
is so on the terms $E^{pq}_2$, and this proves the
bijectivity of~$(*)$.

\begin{corollary}
\label{XII.4.5}
The functor which associates $X^{\an}$ to every proper $\CC$-scheme
$X$ is fully faithful.
\end{corollary}

We must show that, if $X$ and $Y$ are two proper
$\CC$-schemes, the canonical map
$$
\Hom_{\CC}(X,Y)\to\Hom(X^{\an},Y^{\an})
$$
is bijective. Now, to give a morphism from~$X$ to~$Y$ (\resp from
$X^{\an}$ to $Y^{\an}$) is equivalent to giving its graph,
\ie a closed subscheme $Z$ of~$X\times Y$ (\resp a closed
analytic subspace $\othercal{Z}$ of~$X^{\an}\times
Y^{\an}$), such that the restriction of the first projection
$X\times Y\to X$ to~$Z$ (\resp of~$X^{\an}\times Y^{\an}\to
X^{\an}$ to~$\othercal{Z}$) is
% original p. 332
an isomorphism. As the datum of a closed subscheme
of~$X\times Y$ (\resp of a closed analytic subspace
of~$X^{\an}\times Y^{\an}$) is equivalent to that of a coherent
sheaf of ideals on~$\cal{O}_{X\times Y}$ (\resp on~$\cal{O}_{X^{\an}\times Y^{\an}}$), the corollary follows
from~\Ref{XII.4.4}.

\begin{corollary}
\label{XII.4.6}
Let $X$ be a proper $\CC$-scheme. The functor which associates $X'^{\an}$ to every
finite (\resp finite étale) scheme $X'$ over~$X$
is an equivalence of the category of finite
(\resp finite étale) schemes over~$X$ with the category
of finite (\resp finite étale) analytic spaces over~$X^{\an}$.
\end{corollary}

Indeed, to give a finite morphism $X'\to X$ (\resp $X'^{\an}\to
X^{\an}$) is equivalent to giving a coherent sheaf
of algebras on~$\cal{O}_X$ (\resp on~$\cal{O}_{X^{\an}}$)
\cite[\no 19~\S 5~th\ptbl 2]{XII.4}. The corollary therefore follows from~\Ref{XII.4.4}
in the case without ``resp.'', and the ``resp.''\ case is deduced from it taking
into account~\Ref{XII.3.1}~(iii).
